\section{迭代训练管线}
\subsection{Self-Rewarding 结构拆解}
自我奖赏闭环包含三个角色：Task Generator 负责从 Self-Instruct seed 派生任务、Solution Generator 给出具 Chain-of-Thought 的初始回答、Reward Model 执行 cross-check 并计算 reward。Weston 认为，以 Self-Instruct 产生的任务更接近实际需求（代码、数学、评测），因此这一步不要简单地重复已有的 instruction dataset。

\begin{figure}[H]
\centering
\includegraphics[width=0.8\textwidth]{frame-01-selfreward.jpg}
\caption{讲者在早期章节展示 self-improving 概念图，体现 three-role pipeline。\protect\footnotemark}
\end{figure}
\footnotetext{视频画面时间区间：00:02:35--00:02:37。}

\subsection{Self-Instruct 与交叉验证}
为了保持任务广度，pipeline 使用 Self-Instruct 批量生成新的 prompts。每轮回答后，模型生成 verified response，再与初始 draft 配对。如果 meta judge 判定两者不同，则用 verified response 更新 reward 模型。这个过程类似 ``Self-Feeding Chatbot'' 的对话验证，但更专注 reasoning chain，而不是简单的对话回复。

\begin{importantbox}{Self-Instruct 与 verified response 的协同}
\begin{itemize}
    \item Self-Instruct：用少量 seed prompt 自动翻倍出数学、代码和判断型任务。
    \item Verified response：模型判断自己是否错，即便原始 draft 把答案写出来，也必须生成新的、经过验证的版本。
\end{itemize}
\end{importantbox}

\begin{figure}[H]
\centering
\includegraphics[width=0.78\textwidth]{frame-02-iterative.jpg}
\caption{Self-Rewarding pipeline 图示，强调 verified response 的迭代更新。\protect\footnotemark}
\end{figure}
\footnotetext{视频画面时间区间：00:36:00--00:36:02。}

\begin{figure}[H]
\centering
\includegraphics[width=0.75\textwidth]{frame-03-selfinstr.jpg}
\caption{Self-Instruct 中任务生成与验证回答的流程。\protect\footnotemark}
\end{figure}
\footnotetext{视频画面时间区间：00:44:40--00:44:42。}

\subsection{训练数据与微调策略}
Weston 在讲座中反复提到 `Self-Feeding Chatbot'、`SelfT evaluators' 以及 `verified response iterations'，这些组件构成了训练数据的主线。每一轮 self-rewarding 都会加入 alr generation（instruction, question, answer, reasoning），然后在 meta judge 指示下筛选高质量回答。值得关注的是，在实际实验中，用三轮 self-rewarding 训练后的模型与 verified response 的一致性大幅提升，正如 Weston 所说 ``\textit{I presently love the simplicity of it that we didn't have to do this large amount of Chain of Thought where it sort of overthinks itself}''。

\begin{knowledgebox}{Self-Instruct 的采样策略}
Self-Instruct 以少量人工 seeds 为起点，拟监督地生成 instruction → input → output 的三元组。Weston 推荐每轮覆盖语言、数学、编程、逻辑判断四类任务，以防训练数据集中只偏向某一种 reasoning。
\end{knowledgebox}

\subsection{本章小结}
Self-Rewarding 训练用 pipeline 将 Task Generator、Solution Generator、Reward Model 连成一个闭环。Self-Instruct 保证任务多样、Verified Response 让 reward 模型关注错误模式，而 Self-Feeding/Meta judge 进一步验证结果是否值得采纳。

